% !PRINTED 161
\weberpage{161}

Dies ergiebt sich einfach daraus, dass eine Transposition, z.\,B. $(1, 2)$, der Substitution
\[
\begin{pmatrix} 0, & 1, & 0, & \dots & 0 \\ 1, & 0, & 0, & \dots & 0 \\ 0, & 0, & 1, & \dots & 0 \\ \cdot & \cdot & \cdot & \cdot & \cdot \\ 0, & 0, & 0, & \dots & 1 \end{pmatrix},
\]
deren Determinante $-1$ ist, entspricht, und dass man alle Permutationen der ersten Art aus einer geraden, und alle Permutationen der zweiten Art aus einer ungeraden Anzahl von Transpositionen zusammensetzen kann.

\begin{center}
\S.\,40.

\gesperrt{Die Invarianten von endlichen Gruppen linearer Substitutionen.}
\end{center}

Wir beschränken jetzt unsere Betrachtungen auf \gesperrt{endliche Gruppen} linearer homogener Substitutionen. Das Beispiel der Permutationsgruppen zeigt, dass es solche Gruppen giebt. Andere werden wir später noch kennen lernen.

Wir bezeichnen eine solche Gruppe mit $S$, und ihre Substitutionen mit $A, B, C, \dots$ Werden in irgend einer Function der Variablen $x_1, x_2, \dots, x_n$ die Variablen $(x)$ durch $A(x)$ ersetzt, so sagen wir, \gesperrt{die Substitution $A$ werde auf die Variablen $(x)$ angewandt}. Wir definiren nun zunächst folgenden Begriff:

1. \gesperrt{Eine Form $\Phi(x_1, x_2, \dots, x_n)$ von $n$ Variablen heisst eine Invariante oder invariante Form der Gruppe $S$, wenn sie ungeändert bleibt, wenn auf die Variablen $(x)$ die sämmtlichen Substitutionen $A, B, C, \dots$ der Gruppe $S$ angewandt werden.}

Wie früher (Bd.\,I, \S.\,15) ist hier unter einer Form eine ganze homogene Function der Variablen $x$ zu verstehen.

Wir können diese Definition auch durch die Formeln:
\weq{1}{\Phi[A(x)] = \Phi[B(x)] = \Phi[C(x)] \dots}
ausdrücken.

Dass es für jede endliche Gruppe $S$ invariante Formen in beliebiger Menge giebt, ist leicht einzusehen. Man braucht nur
%CONT>>
